% Fichier complété par tools/complete_missing_corrections.py.
% Source : ELEC/ELEC-04-MCC/1023_MCC/1023_MCC.tex
% Statut : rédigé (10 question(s)).
\begin{corrigebox}[Corrigé rédigé]
\CorrigeQuestion{1}
Schéma : alimentation continue de l'inducteur séparé et alimentation
            commandée de l'induit. L'induit est modélisé par $R_a$, $L_a$ et la f.é.m.
            opposée $E=K\Omega$; le couple est $T_{em}=KI$.
\CorrigeQuestion{2}
À l'arrêt, $\Omega=0$, donc $\boxed{E=0}$.
\CorrigeQuestion{3}
Le modèle équivalent est une source $U$ alimentant en série $R_a$, $L_a$
            et la source de tension $E$ opposée au courant $I$.
\CorrigeQuestion{4}
$U=E+R_aI+L_a\,dI/dt$. En régime établi au démarrage, $E=0$ et
            $U_d=R_aI_d=\boxed{1,2R_aI_N}$. Si l'inductance est prise en compte, une rampe
            de tension limite $dI/dt$.
\CorrigeQuestion{5}
On utilise un hacheur quatre quadrants ou un pont redresseur commandé,
            piloté par une boucle de courant interne et une boucle de vitesse externe.
\CorrigeQuestion{6}
$P_a=UI$ pour l'induit; on remplace $U$ et $I$ par les valeurs nominales
            de la plaque.
\CorrigeQuestion{7}
$P_{abs}=UI+U_fI_f$ pour une excitation indépendante.
\CorrigeQuestion{8}
$P_J=R_aI^2+R_fI_f^2$.
\CorrigeQuestion{9}
$P_u=P_{abs}-P_J-P_{autres}$ et
            $\eta=P_u/P_{abs}$. Ici $P_{autres}=27\,\mathrm{kW}$.
\CorrigeQuestion{10}
$T_u=P_u/\Omega_N$, avec $\Omega_N=2\pi n_N/60$.
            $T_{em}=P_{em}/\Omega_N=EI/\Omega_N=KI$.
\end{corrigebox}
